%%
%% fall-lecture11.tex
%% 
%% Made by Alex Nelson <pqnelson@gmail.com>
%% Login   <alex@lisp>
%% 
%% Started on  2026-10-22T11:28:51-0700
%% Last update 2026-10-22T11:28:51-0700
%% 

\lecture{}

\begin{node}
We will study a number of examples of incompressible viscuous fluids
in 2-dimensions.
\end{node}

\subsection{Incompressible 2-dimensional Navier--Stokes equations}

\begin{assume}
Parallel viscuous flow (so the velocity vector looks like $\vec{u}=u\widehat{x}$).
\end{assume}

\begin{assume}
Viscosity is constant $\mu=\mbox{const.}$
\end{assume}

\begin{node}
The continuity equation becomes
\begin{equation}
\partial_{x}u+\partial_{y}v=0.
\end{equation}
\end{node}

\begin{node}[Momentum equations]
The momentum equations become
\begin{subequations}
  \begin{align}
\partial_{t}u+u\partial_{x}u+v\partial_{y}u &= \frac{-1}{\rho}\partial_{x}p+\nu(\partial_{x}^{2}u+\partial_{y}^{2}u)\\
\partial_{t}v+u\partial_{x}v+v\partial_{y}v &= \frac{-1}{\rho}\partial_{x}p+\nu(\partial_{x}^{2}v+\partial_{y}^{2}v)
  \end{align}
\end{subequations}
where $\nu := \mu/\rho$ is called the \define{Kinematical Viscosity}.

Observe: the continuity equation plus the momentum equations are 3
equations in 3 unknowns.
\end{node}

\begin{node}
Under the assumptions of parallel viscous flow: $v=0$, $u=u(y,t)$, and
we have the equations of motion become
\begin{subequations}
  \begin{align}
\partial_{t}u &= \frac{-1}{\rho}\partial_{x}p+\nu\partial_{y}^{2}u\\
0 &=\frac{-1}{\rho}\partial_{y}p.
  \end{align}
\end{subequations}
The $\partial_{t}v=\dots$ equation for conservation of linear momentum
reduces to the tautology $0=0$.
\end{node}

\subsection{Rayleigh Problem}

\begin{problem}
We have one infinitely long wall located at $y=0$. The fluid occupies
some region $y>0$. The wall is stationary at time $t=0$. At time
$t>0$, the wall's velocity becomes constant $\vec{u}_{\text{wall}}=u_{0}\widehat{x}$
moving in the $x$-direction. (The intuition is that the fluid near the
wall slowly accelerates until it reaches velocity
$\vec{u}=\vec{u}_{\text{wall}}$, and a boundary layer propagates in
the fluid moving in the $x$-direction.)
\end{problem}

\begin{node}
There is no length scale in this problem. (We need at least two
distinct points to form a length scale.)
\end{node}

\begin{node}
There is no time scale in this problem.
\end{node}

\begin{node}[Boundary conditions]
We have the two boundary conditions given by the ``no-slip''
condition:
\begin{equation}
u(y=0,t)=u_{\text{wall}}.
\end{equation}
We also have the condition that $\vec{u}$ vanishes ``at infinity'':
\begin{equation}
\lim_{y\to\infty}u(y,t)=0.
\end{equation}
\end{node}

\begin{node}[Similarity solution]
Since there is no time or length scale, we can form a self-similar
solution. This means we use the ans\"{a}tz
\begin{equation}
u(y,t)=u_{\text{wall}}f(\eta)
\end{equation}
where $f$ is some function to be determined (solving our differential
equation), and
\begin{equation}
\eta=\frac{y}{2\sqrt{\nu t}}.
\end{equation}
Observe $\eta$ is dimensionless.

Now, the momentum equation gives us
\begin{equation}
\partial_{t}(u_{\text{wall}}f(\eta))=\nu\partial_{y}(u_{\text{wall}}f(\eta)),
\end{equation}
since $\partial_{x}p=0$ in this case. Then we obtain
\begin{equation}
f''(\eta)+2\eta f'(\eta)=0.
\end{equation}
Let $g(\eta)=f'(\eta)$, then this is the first-order differential
equation $g'(\eta)+2\eta g(\eta)=0$, or
\begin{equation}
\ln(g(\eta)/g(0))=-\eta^{2}\implies f'(\eta)=f'(0)\exp(-\eta^{2}),
\end{equation}
which has general solution
\begin{equation}
f(\eta)=C_{2}+\frac{\sqrt{\pi}C_{1}}{2}\erf(\eta)
\end{equation}
where $\erf(z)=(2/\sqrt{\pi})\int^{z}_{0}\exp(-x^{2})\,\D x$ is the
error function.
Then by the $u(y=0)=u_{\text{wall}}$ boundary condition gives us
$f(0)=1$ (i.e., $C_{2}=1$). Then $u(y=\infty)=0$ boundary condition
gives us $f(\eta)=1-\erf(\eta)=\erfc(\eta)$:
\begin{equation}
u(y,t)=u_{\text{wall}}\erfc\left(\frac{y}{2\sqrt{\nu t}}\right).
\end{equation}
\end{node}

\subsection{Couette--Poiseuille Problem}

\begin{problem}[Couette problem]
We have 2 parallel, infinitely long walls separated by a distance $h$.
We choose coordinates such that the walls lie along [bottom] $y=0$ and
[top] $y=h$. This is assumed to be a steady problem (i.e., $u=u(y)$)
and the bottom wall is stationary and the top wall moves at constant
velocity $U$.
\end{problem}

\begin{problem}[Poiseuille]
We also know pressure gradient is constant $\partial_{x}p=\mbox{const}$.
\end{problem}

\begin{node}
The continuity equation
\begin{equation}
\partial_{x}u+\partial_{y}v=0.
\end{equation}
This implies $u=u(y)$ is a function of $y$ alone.
\end{node}

\begin{node}
The $y$-component of the momentum conservation equation gives us
\begin{equation}
0=\frac{-1}{\rho}\partial_{y}p.
\end{equation}
This implies
\begin{equation}
\partial_{y}p=0,
\end{equation}
which has its solution be
\begin{equation}
p=p(x).
\end{equation}
That is to say, pressure is a function of $x$ alone.
\end{node}

\begin{node}
The $x$-component of the conservation of linear momentum equation is
\begin{equation}
\partial_{t}u+v\partial_{y}u=\frac{-1}{\rho}\partial_{x}p+\nu(\partial_{x}^{2}u+\partial_{y}^{2}u).
\end{equation}
But the only nonzero terms surviving gives us
\begin{equation}
0=\frac{-1}{\rho}\partial_{x}p+\nu\partial_{y}^{2}u.
\end{equation}
Then we have
\begin{equation}
\frac{1}{\rho}\partial_{x}p=\nu\partial_{y}^{2}u.
\end{equation}
\end{node}

\begin{node}[Poiseuille problem]
When $\partial_{x}p$ is constant and we have boundary conditions
\begin{equation}
u(y=0)=0,\quad\mbox{and}\quad u(y=h)=0.
\end{equation}
Then our differential equation becomes
\begin{equation}
\mbox{(constant)}=\nu\partial_{y}^{2}u.
\end{equation}
The general solution looks like
\begin{equation}
u = \frac{1}{2\mu}(\partial_{x}p)y^{2}+A_{1}y+A_{2},
\end{equation}
where $A_{1}$ and $A_{2}$ are constants to be determined. We see
\begin{equation}
u(y=0)=A_{2}=0\implies A_{2}=0.
\end{equation}
Similarly,
\begin{equation}
u(h)=\frac{1}{2\mu}(\partial_{x}p)h^{2}+A_{1}h=0,
\end{equation}
which implies
\begin{equation}
A_{1}=\frac{-1}{2\mu h}(\partial_{x}p)h^{2}=\frac{-1}{2\mu}(\partial_{x}p)h.
\end{equation}
Therefore
\begin{equation}
u(y) = \frac{1}{2\mu}(\partial_{x}p)(y^{2}-yh).
\end{equation}
The velocity profile is a parabola in the Poiseuille problem.
\end{node}

\begin{node}[Couette flow]
At the other extreme, when $\partial_{x}p=0$, we see that
\begin{equation}
\partial_{y}^{2}u=0\implies u=A_{1}y+A_{2}.
\end{equation}
Now we impose the boundary conditions $u(0)=A_{2}=0$ and
\begin{equation}
u(h)=A_{1}h=U\implies A_{1}=\frac{U}{h}.
\end{equation}
This means the general solution to the Couette flow looks like
\begin{equation}
u(y)=U\frac{y}{h}.
\end{equation}
\end{node}

\begin{remark}
A few things to think about and compute:
\begin{enumerate}
\item $\tau_{12}|_{\text{wall}}$ tells us the shear stress of the
  wall, it tells us the wall drag. It's also called the ``skin stress''.
  Explicitly, $\tau_{12}|_{\text{wall}}=\mu\partial_{y}u|_{\text{wall}}$
\item The dissipation function $\Phi:=\tau_{ij}\partial_{j}u_{i}$ may
  be computed in this situation as $\Phi=\mu(\partial_{y}u)^{2}$. 
\end{enumerate}
\end{remark}

\subsection{Oscillating boundary layer}

(This is also called the ``Stokes boundary layer problem'')

\begin{problem}
We have one infinitely long wall located at $y=0$. The wall moves at
velocity
\begin{equation}
u_{\text{wall}}=u_{0}\cos(\omega t),
\end{equation}
where $\omega\neq0$ is some constant. Assume there is no pressure
gradient in this problem. Solve for $u(y,t)$.
\end{problem}

\begin{node}
The equations of motion (specifically the conservation of linear
momentum in the $x$-direction) is:
\begin{equation}
\partial_{x}u + 0 = 0 + \nu(0 + \partial_{y}^{2}u).
\end{equation}
We have the boundary conditions
\begin{subequations}
\begin{equation}
u(y=0,t)=u_{\text{wall}}=u_{0}\cos(\omega t),
\end{equation}
and
\begin{equation}
\lim_{y\to\infty}u(y,t)=0.
\end{equation}
\end{subequations}
\end{node}

\begin{node}
We know by linearity that
\begin{equation}
u = \Re\left(\E^{\I \omega t}f(y)\right),
\end{equation}
where $f(-)$ is ``some function'', is the general form of the solution.
\end{node}

\begin{node}
We can also use Fourier analysis, where $f(y)$ here is the Fourier
mode at the particular frequency. If we apply the Fourier transform
for time, then the partial differential equation becomes
\begin{equation}
\I\omega\widehat{u}=\nu\partial_{y}^{2}\widehat{u}.
\end{equation}
This has its solution be
\begin{equation}
\widehat{u}(y)=A\exp\left(-(1+\I)\frac{y}{\delta}\right)+B\exp\left((1+\I)\frac{y}{\delta}\right)
\end{equation}
where $\delta=\sqrt{2\nu/\omega}$, and $A$ and $B$ are constants to be
determined (by the boundary conditions).
\end{node}

\begin{node}
We can perform the inverse Fourier transform, then impose the boundary
conditions, which will give us the solution:
\begin{equation}
u(y,t)=u_{0}\E^{-y/\delta}\cos\left(\omega t-\frac{y}{\delta}\right).
\end{equation}
\end{node}

\begin{remark}
The factor $\E^{-y/\delta}$ tells us the how the magnitude of the
oscillation decays as we move away from the wayy. The $\delta$ gives
us an effective boundary length (i.e., a length scale).
\end{remark}

\begin{remark}
We also have a phase $-y/\delta$ in the cosine of the solution, which
gets more noticeable as we move away from the wall.
\end{remark}

\begin{remark}
Compare this to the Yayleigh problem when the boundary layer
(where $0<u(y,t)<u_{\text{wall}}$) propagates to infinity. But here,
the boundary layer does not.

Also the Rayleigh problem had
\begin{equation}
\frac{u}{y}\sim\nu\frac{u}{\delta_{R}^{2}}
\end{equation}
where $\delta_{R}\sim\sqrt{\nu t}$. But here with the oscillating
wall, we have instead
\begin{equation}
\omega u\sim\nu\frac{u}{\delta^{2}},
\end{equation}
where $\delta=\sqrt{2\nu/\omega}$.
\end{remark}

\subsection{Stokes flow (low Reynolds number)}

\begin{problem}
The Stokes flow assumes $u_{j}\partial_{j}u_{i}\ll\nu\partial_{j}^{2}u_{i}$,
so we can ignore the small nonlinear terms. The momentum equations
become a linear system of partial differential equations.
\end{problem}

\begin{node}
Further, people tend to assume steady flow or $\partial_{t}u$
contributions are sufficiently small.

This generally makes the left-hand side of the conservation of linear
momentum identically zero.
\end{node}
